% Осень26!A13; нумерация внутри столбца без пустых строк.
% Источник: лекция 5, §10 (строка 347)
\ticket{11}{Лемма об оценке квадратичной формы и об эквивалентных нормах}
\label{ticket:11}

\begin{definition}
Нормы $\rho_1,\rho_2$ на одном пространстве называются \textbf{эквивалентными}, если существуют $c_1,c_2>0$ такие, что
\[
  c_1\rho_1(x)\le\rho_2(x)\le c_2\rho_1(x)
  \qquad\text{для всех }x.
\]
\end{definition}

\begin{lemma}
\begin{enumerate}
  \item Для положительно определённой квадратичной формы $Q$ на $\mathbb R^m$ существует $\gamma>0$ такое, что $Q(x)\ge\gamma|x|^2$ для всех $x$.
  \item Всякая норма $\rho$ на $\mathbb R^m$ эквивалентна евклидовой: существуют $c_1,c_2>0$ такие, что $c_1|x|\le\rho(x)\le c_2|x|$.
\end{enumerate}
\end{lemma}

\begin{proof}
На единичной сфере $Q$ непрерывна и положительна.
По компактности $\gamma:=\min_{|u|=1}Q(u)>0$.
Для $x\ne0$
\[
  Q(x)=|x|^2Q(x/|x|)\ge\gamma|x|^2.
\]
При $x=0$ оценка очевидна.

Для произвольной нормы $\rho$ по неравенствам треугольника и Коши--Буняковского
\[
  |\rho(x)-\rho(y)|\le\rho(x-y)
  \le\sum_{j=1}^m|x_j-y_j|\rho(e_j)
  \le C|x-y|,
\]
где $C=(\sum_j\rho(e_j)^2)^{1/2}$.
Следовательно, $\rho$ непрерывна, и на единичной сфере достигаются
\[
  c_1:=\min_{|u|=1}\rho(u)>0,
  \qquad c_2:=\max_{|u|=1}\rho(u)<\infty.
\]
Из $\rho(x)=|x|\rho(x/|x|)$ при $x\ne0$ следует требуемая двусторонняя оценка.
\end{proof}
